\chapter[Sermon 94. Yet Again Is Described the Seat of the Everlasting Country in Heaven.]{Yet Again Is Described the Seat of the Everlasting Country in Heaven.}
\label{sermon:094}
\markright{Sermon 94. Yet Again Is Described the Seat of the Everlasting ...}

And he who spoke with me had a golden reed to measure the city, and its gates, and its walls. And the city was built squarely, and its length was equal to its breadth. And he measured the city with the reed twelve thousand six hundred furlongs, and the length, breadth, and height of it were equal. And he measured the wall of it one hundred forty-four cubits, according to the measure of a man which the angel had. And the construction of the wall was of jasper. And the city was of pure gold, like clear glass, and the foundations of the walls and of the city were adorned with all manner of precious stones. The first foundation was jasper, the second sapphire, the third chalcedony, the fourth emerald, the fifth sardonyx, the sixth sardius, the seventh chrysolite, the eighth beryl, the ninth topaz, the tenth chrysoprase, the eleventh jacinth, the twelfth amethyst. And the twelve gates were twelve pearls, and each gate was one pearl, and the street of the city was pure gold, as a thoroughfare of shining glass.

He proceeds in describing the blessed dwellings and that life of the world to come, under the image of a most beautiful and most excellent city. We must understand all things not according to the literal meaning, but according to the spirit. All things are said for our comfort, and so that we should boldly despise this world and its pleasures, and the fury of persecutors; and should always desire those great and everlasting good things promised to us, as we have heard in the description. Indeed, we have seen four particular things of this heavenly city as if in a lively picture: what light it has, what walls, what gates, and foundations. Now, in the fifteenth place, follows what the width, capacity, or largeness of this city is. For cities are commended because of their size. It is necessary that the greatest number of citizens should have the largest city.

Therefore comes forth a messenger of this city, an angel sent to John from heaven, holding in his hand a reed, that is a long pole or measuring rod, not of wood or lead, but of gold. By the measuring, he would have us understand the quantity of the blessed seat. In the meter and in the measure, we shall not need to seek any great mysteries. For the eternal wisdom and providence of God has prepared seats for his chosen; and that in a golden order, that is to wit, most purified, which is signified by the golden reed or measure. For the judge in Matthew provokes the sheep to take the inheritance, prepared from the beginning of the world. He alone knows also, who are his.

The situation of the city is declared, planted in a square, whereby is signified the strength and stability of the blessed in heaven. For the place is no ball, bowl, or globe, rolling and easy to turn. Neither need we to doubt of the certainty thereof. For hope shames no man, and he that believes in Christ shall never be confounded.

Moreover, the length, breadth, and height of the city are equal. Every side, in its square, has twelve thousand furlongs, which make in all forty-eight thousand in the whole circuit. Concerning the furlong, what and how much it contains, I see that learned men vary. Pliny in book 2, chapter 23, attributes to a furlong one hundred and twenty-five paces, that is to say, six hundred and twenty-five feet. If you now calculate these things and divide them into miles, you will find that the city is most ample and large. Some reckon it one hundred and fifty German miles. Hereby I suppose to be signified that the place and space is great enough, whatever innumerable multitudes of angels, of blessed spirits, and men shall flit into the blessed seat and dwell therein. As also the Lord in the Gospel said: “In my Father’s house are many mansions.” In chapter 30 of Isaiah toward the end, it is shown that there shall be space and place enough in hell also for the wicked. And the equality on every side declares that men of people or nations shall have no prerogative. For whether you be of the East or of the West, whether you be Greek or Barbarian, so that you be faithful, you shall be received by the Lord. Moreover, in the gospel equality is declared, while the penny is paid not only to him who wrought in the vineyard all day long or half the day, but unto him also who came into the vineyard in the evening.

Note 94.6: The height of the wall is undoubtedly immeasurable. From this we gather that the blessedness is most certain, and that no one can enter it except through the gates. For no one can climb over such a height, nor scale those walls, whether he be an enemy who would trouble them, or a hypocrite who attempts to gain heaven by stealth.

\index{Resurrection}Where he says, “and he measured the wall thereof, an hundred and forty-four cubits,” it cannot certainly agree with furlongs; therefore, we must understand it of the thickness of the wall. By this again is figured the strength and security of the blessed. It is added how the Angel measured with the measure of a man, which the Angel had: that is to say, that the Angel measured the customary cubits and furlongs used by men. Therefore, this Angel had the same measure in this measuring as is commonly used by men. For so would he signify that the place of eternal felicity should be determinate and certain. For there shall be after the resurrection true and determinate bodies. If there be any other mystery herein, perhaps it is the same which the Lord spoke of in Luke, namely of the blessedness of the faithful in another world: they are like angels, and are the children of God, since they are the children of the resurrection. If any man will account these numbers more exactly and show higher mysteries, I will gladly give place. I suppose here rather celestial things to be figured than either arithmetical numbers or geometrical proportions to be taught. Nevertheless, I can willingly grant that those arts help to the understanding of the Scriptures.

\index{Rome}In the sixth place, he discusses the matter of this heavenly City. For cities are praised for their substance and material. The saying of Caesar Augustus is well known, who is said to have spoken of Rome: “I found it of brick, I leave it of marble.” And cities built of stone are justly preferred to those that are of timber, and those built of squared free stone to those made of rough stone. But what is the building or material of the celestial City? He declares that by five parts or members. First, the walls are of jasper. Let no man here imagine carnal things. The jasper is green. The celestial City always flourishes; God’s protection never fails.

2. The city itself, that is to say, the buildings in the city, the palaces and houses, are pure gold. For all things are purified in the eternal realm. There is no uncleanness, no evil affections, there shall be no trouble or pain. As the Lord said also in the 19th Chapter of Matthew, disputing against the Sadducees. Therefore, just as gold is most tried and pure, so shall the heavenly dwelling be most clean. Therefore, the bodies that shall dwell in heaven must also be clarified or glorified. He adds that this most pure gold is not glass, but in brightness represents most pure and shining glass. For in heaven all things are clear. There we shall be seen face to face. There we shall most perfectly know all things.

3. And first he says generally that the foundations of the city are beautified with all manner of precious stones; after which he recites particularly the stones that are most excellent. Doubtless nothing is more precious, nothing more excellent, than Christ, the foundation of our salvation, than the apostolic doctrine, whereby we are induced to the knowledge of Christ and of our salvation. And he sets in order twelve stones, to the intent we should understand that there is not one precious stone alone placed for the foundation, but a row of one sort in such a length as the side is square, and so consequently likewise in all parts of the square. For the first order therefore is placed a jasper stone, that is to say, in the first place of the foundation, jasper stones are set in their rank; again in the next row upon the jaspers are laid sapphires, throughout the whole space, in such length as the foundation was, and so consequently the other stones were couched and laid in order. By all which is signified that the foundation of our salvation is both most excellent and sure. Which we ought of right to set more by than by the price of all the jewels in the earth. And there are found men godly and beneficial, who, bestowing or selling these earthly jewels, according to the Apostle’s doctrine in 1 Timothy 6, prepare for themselves a good foundation in another world. There are foolish people who are overly in love with jewels, and many times instead of precious stones that cost very much when polished, they obtain glass. Truly worthy doubtless to be deceived. Verily, precious stones have their use and virtues, neither were they made by God in vain. But we must always remember that saying of the wise man: all things are not suitable for all men.

\index{Apocalypse (Book of)}\index{Jerome, Saint}\index{Epiphanius of Salamis}\index{Arethas of Caesarea}\index{Aquinas, Thomas}By the listing of precious stones, he appears to have alluded to the precious stones set in the attire of the high priest, as described in Exodus 28. I do not doubt that John took these things partly from Isaiah 54, which Jerome expounding, directs those who desire to know more about stones to Epiphanius and to Pliny’s 37th book on Natural History. Arethas in his commentaries applies the twelve precious stones to the twelve Apostles of Christ. Furthermore, the writings of Bede on this passage remain; from whom Thomas Aquinas took what he has in his commentaries on the Apocalypse. I see no way that I can profitably linger longer in this treatise. Therefore, I refer the curious reader to these authors; it is enough for me to have shown that these costly jewels signify the excellence of the foundation of our health and salvation.

Moreover, in the fourth place, the matter of the gates is declared. They were all of one whole pearl, each of them, the price of which is exceedingly great. The gate of heaven is Christ, and the gatekeepers of heaven are the Apostles, as has been declared before. Therefore, the gates are most precious and most strong. In the thirteenth chapter of the Gospel according to Matthew, Christ himself and the salvation that comes from him are compared to a pearl, for which the merchant sells all that he has, thinking himself rich enough if he may obtain this pearl.

5. In the fifth place, the street is also described, what it is. In earthly cities, the streets are often muddy, even if the cities themselves are otherwise famous and noble. Where they are noteworthy, they are paved with stone or brick; but the street of our city is paved with pure and bright gold. For in heaven there is no unpleasantness, no obscure darkness. All these things are undoubtedly spoken beautifully; yet far greater things must be understood and imagined. We must strive with all our might that the thing which the tongue of man cannot utter, nor our mind here conceive due to its greatness and excellence, we may at length behold in heaven presently, and experience it in our glorified bodies through Jesus Christ our Lord.
